\documentclass[12pt,a4paper,oneside]{book}

% ============================================================
% Paket dasar
% ============================================================
\usepackage[utf8]{inputenc}
\usepackage[T1]{fontenc}
\usepackage[indonesian]{babel}
\usepackage{mathptmx}
\usepackage[scaled=0.92]{helvet}
\usepackage{geometry}
\geometry{left=3cm,right=2cm,top=2cm,bottom=2cm}
\setlength{\headheight}{15pt}
\usepackage{graphicx}
\usepackage{float}
\usepackage{listings}
\usepackage{xcolor}
\usepackage{fancyhdr}
\usepackage{tocloft}
\usepackage{setspace}
\usepackage{hyperref}
\usepackage{booktabs}
\usepackage{longtable}
\usepackage{array}
\usepackage{tabularx}
\usepackage{enumitem}
\usepackage{caption}
\usepackage{titlesec}
\usepackage{microtype}
\usepackage{amsmath}
\usepackage{tikz}
\usetikzlibrary{arrows.meta,positioning,shapes.geometric}
\usepackage{ragged2e}
\usepackage{etoolbox}
\usepackage{textcomp}
\usepackage{multirow}

% ============================================================
% Pengaturan umum
% ============================================================
\onehalfspacing
\setlength{\parindent}{1cm}
\setlength{\parskip}{0.15cm}
\setcounter{tocdepth}{2}
\setcounter{secnumdepth}{3}
\renewcommand{\arraystretch}{1.15}
\setlength{\emergencystretch}{2em}
\sloppy

\captionsetup{font=small,labelfont=bf,justification=centering}
\captionsetup[table]{skip=6pt}

\hypersetup{
  colorlinks=true,
  linkcolor=black,
  citecolor=black,
  urlcolor=blue!50!black,
  pdftitle={Laporan UTS Pengembangan Perangkat Lunak - Buku Servis},
  pdfauthor={Rizendy Affarin},
  pdfsubject={Laravel, Vue, Docker, FrankenPHP, CI/CD},
  bookmarksnumbered=true
}

% ============================================================
% Header dan footer
% ============================================================
\pagestyle{fancy}
\fancyhf{}
\fancyhead[L]{\nouppercase{\leftmark}}
\fancyhead[R]{\thepage}
\renewcommand{\headrulewidth}{0.4pt}
\fancypagestyle{plain}{%
  \fancyhf{}%
  \fancyhead[R]{\thepage}%
  \renewcommand{\headrulewidth}{0pt}%
}

% ============================================================
% Judul bab dan daftar
% ============================================================
\titleformat{\chapter}[display]
  {\normalfont\Large\bfseries}
  {\chaptertitlename\ \thechapter}{0.5ex}{\Large}
\titlespacing*{\chapter}{0pt}{-0.2cm}{0.8cm}

\renewcommand{\cftchapfont}{\bfseries}
\renewcommand{\cftchappagefont}{\bfseries}
\renewcommand{\cftfigpresnum}{Gambar~}
\renewcommand{\cfttabpresnum}{Tabel~}
\setlength{\cftfignumwidth}{2.5cm}
\setlength{\cfttabnumwidth}{2.2cm}

% ============================================================
% Gaya kode
% ============================================================
\definecolor{codegreen}{RGB}{28,110,70}
\definecolor{codeblue}{RGB}{24,70,140}
\definecolor{codered}{RGB}{150,45,45}
\definecolor{codegray}{RGB}{110,110,110}
\definecolor{codebg}{RGB}{247,247,244}

\lstdefinestyle{proyek}{
  backgroundcolor=\color{codebg},
  commentstyle=\color{codegreen}\itshape,
  keywordstyle=\color{codeblue}\bfseries,
  stringstyle=\color{codered},
  numberstyle=\tiny\color{codegray},
  basicstyle=\ttfamily\footnotesize,
  breakatwhitespace=false,
  breaklines=true,
  captionpos=b,
  keepspaces=true,
  numbers=left,
  numbersep=8pt,
  showspaces=false,
  showstringspaces=false,
  showtabs=false,
  tabsize=2,
  frame=single,
  rulecolor=\color{black!25},
  columns=fullflexible,
  upquote=true,
  literate={-}{-}1
}
\lstdefinelanguage{JavaScript}{
  keywords={break, case, catch, const, continue, default, delete, do, else,
            export, extends, finally, for, from, function, if, import, in,
            instanceof, let, new, return, super, switch, this, throw, try,
            typeof, var, void, while, with, yield, async, await, of, class},
  sensitive=true,
  comment=[l]{//},
  morecomment=[s]{/*}{*/},
  morestring=[b]',
  morestring=[b]",
  morestring=[b]`
}

\lstset{style=proyek}
\renewcommand{\lstlistingname}{Kode}
\renewcommand{\lstlistlistingname}{DAFTAR KODE}

% ============================================================
% Perintah bantu
% ============================================================
\newcommand{\repourl}{https://github.com/Ryzen0000001/uts-pl-24-535196-SV-24171}
\newcommand{\ghcrbackend}{ghcr.io/ryzen0000001/uts-pl-24-535196-sv-24171-backend}
\newcommand{\ghcrfrontend}{ghcr.io/ryzen0000001/uts-pl-24-535196-sv-24171-frontend}
\newcommand{\demoakun}{\texttt{demo@example.com} / \texttt{password123}}

\newcommand{\gambarpenting}[3][0.86\textwidth]{%
  \begin{figure}[H]
    \centering
    \includegraphics[width=#1,height=0.78\textheight,keepaspectratio]{#2}
    \caption{#3}
  \end{figure}
}

\newcommand{\gambarlebar}[3][\textwidth]{%
  \begin{figure}[H]
    \centering
    \includegraphics[width=#1]{#2}
    \caption{#3}
  \end{figure}
}

\newcommand{\isipostman}[2]{%
  \IfFileExists{img/#1}{%
    \begin{figure}[H]
      \centering
      \includegraphics[width=0.82\textwidth,height=0.78\textheight,keepaspectratio]{img/#1}
      \caption{#2}
    \end{figure}
  }{%
    \begin{figure}[H]
      \centering
      \fbox{%
        \parbox{0.78\textwidth}{%
          \centering\small
          Tangkapan layar \texttt{#1} belum tersedia saat laporan ini disusun.
          Permintaan dan status yang diuji tetap dicatat pada Tabel~\ref{tab:api}
          dan transkrip Lampiran~\ref{lam:api}.
        }%
      }
      \caption{#2}
    \end{figure}
  }%
}

% ============================================================
% Dokumen
% ============================================================
\begin{document}

% ---------- Halaman judul ----------
\pagenumbering{alph}
\begin{titlepage}
  \thispagestyle{empty}
  \begin{center}
    \includegraphics[width=2.3cm]{img/lambang-ugm.png}\\[0.35cm]
    {\small\bfseries UNIVERSITAS GADJAH MADA}\\[0.1cm]
    {\small SEKOLAH VOKASI}\\[0.1cm]
    {\small PROGRAM STUDI D-IV TEKNOLOGI REKAYASA PERANGKAT LUNAK}\\[1.1cm]

    {\Large\bfseries LAPORAN UJIAN TENGAH SEMESTER}\\[0.25cm]
    {\large Mata Kuliah Pengembangan Perangkat Lunak}\\[1.1cm]

    {\LARGE\bfseries BUKU SERVIS}\\[0.3cm]
    {\Large Aplikasi Pencatatan Servis Motor Berbasis Laravel, Vue, Docker, dan FrankenPHP}\\[1.4cm]

    \begin{tabular}{ll}
      \textbf{Disusun oleh} & : Rizendy Affarin \\
      \textbf{NIM}          & : 24/535196/SV/24171 \\
      \textbf{Kelas}        & : PL5B1 \\
      \textbf{Dosen}        & : Galih Malela Damaraji \\
    \end{tabular}

    \vfill
    {\bfseries YOGYAKARTA}\\
    {\bfseries 2026}
  \end{center}
\end{titlepage}

% ---------- Front matter ----------
\frontmatter
\pagestyle{plain}

\chapter*{ABSTRAK}
\addcontentsline{toc}{chapter}{ABSTRAK}
\phantomsection

Buku Servis adalah aplikasi web untuk mencatat riwayat servis motor. Setiap pengguna
memiliki catatan sendiri yang berisi tanggal servis, odometer, jenis servis, biaya,
bengkel, dan catatan. Laporan ini mendokumentasikan pembangunan aplikasi tersebut
sebagai pemenuhan UTS mata kuliah Pengembangan Perangkat Lunak: REST API Laravel 12
dengan Sanctum, antarmuka Vue 3 yang memakai Axios, MySQL, Docker Compose,
FrankenPHP sebagai server backend, serta pipeline GitHub Actions yang menguji,
membangun, dan mendorong image ke GHCR hanya dari branch \texttt{main}.

Hasil pengujian menunjukkan 16 tes fitur lulus dengan 38 asersi. Seluruh endpoint
selain \texttt{POST /api/login} menolak permintaan tanpa token dengan status 401.
Pembuatan catatan mengembalikan 201, validasi mengembalikan 422, dan operasi baca,
ubah, hapus mengembalikan 200. Ketiga container berjalan sehat; data tetap ada
setelah \texttt{docker compose down} dan \texttt{docker compose up -d}. Backend
terbukti dilayani FrankenPHP 1.13.3 dengan PHP 8.4.27, bukan \texttt{php artisan serve}
atau PHP-FPM. Pipeline CI berjalan pada branch fitur, \texttt{dev}, dan pull request;
job deploy hanya berjalan setelah push ke \texttt{main} dan menghasilkan image
bertag commit SHA di GHCR.

\vspace{0.4cm}
\noindent\textbf{Kata kunci:} Laravel, Sanctum, Vue, Axios, Docker Compose,
FrankenPHP, GitHub Actions, GHCR.

\chapter*{KATA PENGANTAR}
\addcontentsline{toc}{chapter}{KATA PENGANTAR}
\phantomsection

Laporan ini disusun untuk mendokumentasikan pekerjaan UTS Pengembangan Perangkat
Lunak yang berfokus pada integrasi backend, frontend, basis data, container, dan
CI/CD. Aplikasi yang dibangun bernama Buku Servis. Masalah yang diangkat sederhana:
catatan servis motor sering tersebar di buku manual, pesan chat, atau struk yang
mudah hilang, sehingga pemilik motor sulit mengetahui kapan servis terakhir dilakukan
dan berapa biaya yang sudah dikeluarkan.

Saya mengucapkan terima kasih kepada Bapak Galih Malela Damaraji selaku dosen
pengampu, rekan-rekan kelas PL5B1, dan pihak-pihak yang membantu proses pengujian.
Laporan ini berisi kode, tangkapan layar, keluaran perintah, serta kendala yang
benar-benar ditemui selama pengerjaan. Bagian yang belum sempurna saya tulis apa
adanya agar pembaca dapat menilai pekerjaan ini secara jujur.

\vspace{0.6cm}
\begin{flushright}
  Yogyakarta, 24 Oktober 2026\\[1.2cm]
  Rizendy Affarin
\end{flushright}

\begingroup
\singlespacing
\tableofcontents
\listoffigures
\listoftables
\lstlistoflistings
\endgroup

% ---------- Isi ----------
\mainmatter
\pagestyle{fancy}

\chapter{PENDAHULUAN}
\label{chap:pendahuluan}

\section{Latar Belakang}

Motor adalah kendaraan yang perawatannya bergantung pada catatan. Penggantian oli
mesin, kampas rem, ban, aki, dan servis CVT punya interval yang berbeda. Tanpa
catatan, pemilik motor cenderung mengandalkan ingatan atau menebak dari suara mesin.
Akibatnya servis bisa terlambat, biaya membengkak, dan riwayat perawatan tidak bisa
dijelaskan saat motor dijual.

Aplikasi pencatatan sederhana dapat menyelesaikan masalah tersebut asalkan dapat
diakses dari browser, menyimpan data secara permanen, dan membatasi data berdasarkan
pengguna. Karena itu proyek ini tidak hanya berhenti pada CRUD di satu komputer.
Aplikasi harus berjalan sebagai beberapa container yang saling terhubung, memakai
basis data MySQL, memiliki antarmuka Vue, dan melewati pipeline CI/CD. Persyaratan
tersebut memaksa saya berlatih pada hal yang sering terlewat dalam tugas CRUD biasa:
konfigurasi jaringan antar-container, migrasi basis data saat container mulai,
penyimpanan volume, penanganan token, serta pemisahan rahasia dari repositori.

Aplikasi yang dibangun diberi nama \textbf{Buku Servis}. Entitas utamanya adalah
catatan servis dengan field motor, tanggal servis, odometer, jenis servis, biaya,
bengkel, dan catatan. Setiap catatan dimiliki oleh satu pengguna. Autentikasi memakai
Laravel Sanctum dengan token bearer. Registrasi publik tidak dibuka karena ketentuan
tugas menyatakan semua endpoint selain login harus memakai token; akun demo dibuat
oleh seeder saat backend pertama kali berjalan.

\section{Rumusan Masalah}

Berdasarkan latar belakang tersebut, masalah yang diselesaikan dalam laporan ini
adalah:

\begin{enumerate}[leftmargin=1.2cm]
  \item Bagaimana membangun REST API Laravel yang seluruh endpoint selain login
        dilindungi token Sanctum dan mengembalikan status HTTP yang sesuai?
  \item Bagaimana menghubungkan frontend Vue 3 ke API tanpa menulis alamat API di
        dalam komponen?
  \item Bagaimana menjalankan backend, frontend, dan MySQL melalui Docker Compose
        dengan jaringan service-name dan volume permanen?
  \item Bagaimana memastikan backend benar-benar dilayani FrankenPHP, bukan
        \texttt{php artisan serve} atau PHP-FPM?
  \item Bagaimana menyusun alur Git dan GitHub Actions sehingga pengujian berjalan
        di setiap branch, sedangkan deploy hanya terjadi dari \texttt{main}?
\end{enumerate}

\section{Tujuan}

Tujuan yang ingin dicapai:

\begin{enumerate}[leftmargin=1.2cm]
  \item Menyediakan API CRUD catatan servis yang aman dan konsisten.
  \item Menyediakan antarmuka login dan pengelolaan catatan servis yang memakai
        token dari API.
  \item Menjalankan seluruh sistem dengan satu perintah Docker Compose dan
        memastikan data bertahan setelah container dihentikan.
  \item Membuktikan penggunaan FrankenPHP melalui konfigurasi image dan log runtime.
  \item Menerapkan CI/CD, branch protection, serta penyimpanan rahasia di GitHub
        Secrets.
  \item Mendokumentasikan proses, bukti pengujian, dan kendala secara jujur.
\end{enumerate}

\section{Ruang Lingkup}

Ruang lingkup aplikasi dibatasi pada:

\begin{itemize}[leftmargin=1.2cm]
  \item Satu resource CRUD, yaitu \texttt{service\_logs} (catatan servis motor).
  \item Autentikasi token melalui \texttt{POST /api/login}, \texttt{GET /api/me},
        dan \texttt{POST /api/logout}. Tidak ada registrasi publik.
  \item Frontend satu halaman untuk login dan satu halaman untuk daftar, tambah,
        ubah, dan hapus catatan.
  \item Basis data MySQL 8.4 dengan volume bernama.
  \item Backend Laravel 12 + Sanctum yang dilayani FrankenPHP.
  \item CI/CD GitHub Actions dengan job pengujian, build, dan deploy ke GHCR.
\end{itemize}

Di luar ruang lingkup: aplikasi mobile native, notifikasi pengingat servis,
multi-bahasa, pembayaran, dan manajemen bengkel multi-cabang.

\section{Manfaat}

Secara praktis, aplikasi ini memberi cara cepat untuk mencatat servis dan melihat
total biaya perawatan. Secara akademis, proyek ini menjadi latihan integrasi antara
API, SPA, basis data, container, dan pipeline. Dari sisi dokumentasi, laporan ini
menyimpan bukti eksekusi yang dapat diperiksa ulang: keluaran tes, status HTTP,
log FrankenPHP, keluaran \texttt{docker compose ps}, dan tautan commit.

\section{Sistematika Penulisan}

\begin{description}[leftmargin=1.2cm,style=nextline]
  \item[BAB I Pendahuluan] memuat latar belakang, rumusan masalah, tujuan, ruang
        lingkup, manfaat, dan sistematika.
  \item[BAB II Landasan Teori] membahas Laravel, Sanctum, Vue, Axios, Docker,
        FrankenPHP, dan GitHub Actions.
  \item[BAB III Analisis dan Perancangan] menjelaskan kebutuhan, model data,
        rancangan API, rancangan antarmuka, dan rancangan deployment.
  \item[BAB IV Implementasi] memaparkan kode dan konfigurasi yang ditulis.
  \item[BAB V Hasil dan Pembahasan] menyajikan bukti pengujian per CPMK, alur data,
        serta kendala dan solusinya.
  \item[BAB VI Kesimpulan dan Saran] merangkum hasil dan memberi saran lanjutan.
\end{description}

\chapter{LANDASAN TEORI}
\label{chap:teori}

\section{Laravel dan Sanctum}

Laravel adalah framework PHP dengan pola MVC. Pada versi 12, konfigurasi aplikasi
dipusatkan di \texttt{bootstrap/app.php}; routing API dipisahkan dari routing web.
Permintaan API masuk melalui \texttt{public/index.php}, melewati middleware global,
middleware grup, lalu router. Controller mengembalikan respons JSON. Eloquent
menjembatani model dengan tabel MySQL.

Sanctum menyediakan token akses pribadi. Saat login berhasil, server membuat token
dan menyimpan hash-nya di tabel \texttt{personal\_access\_tokens}. Token asli hanya
dikirim sekali ke klien. Pada permintaan berikutnya, klien mengirim header
\texttt{Authorization: Bearer <token>}. Middleware \texttt{auth:sanctum} memeriksa
token tersebut. Jika tidak ada atau tidak valid, permintaan ditolak dengan 401.

\section{Vue 3 dan Axios}

Vue 3 dengan Composition API memakai \texttt{ref}, \texttt{computed}, dan
\texttt{onMounted} untuk mengelola state halaman. Vue Router mengatur navigasi
\texttt{/masuk} dan \texttt{/} serta menjaga halaman daftar agar hanya dapat dibuka
saat token tersedia. Axios dipakai sebagai klien HTTP karena mendukung instance
dengan \texttt{baseURL} dan interceptor.

Interceptor request menambahkan token ke setiap permintaan. Interceptor response
menangani 401 dengan memicu event \texttt{auth:expired}, sehingga aplikasi dapat
kembali ke halaman masuk ketika token dicabut. Alamat API diambil dari
\texttt{import.meta.env.VITE\_API\_URL}; nilai ini diisi saat build, bukan ditulis
di dalam komponen.

\section{Docker dan Docker Compose}

Docker mengemas aplikasi beserta dependensinya ke dalam image. Container adalah
instance image yang berjalan. Docker Compose mendefinisikan beberapa service dalam
satu berkas YAML. Setiap service dapat memiliki image, build context, environment,
port, volume, healthcheck, dan network.

Dua hal penting dalam proyek ini adalah \emph{multi-stage build} dan \emph{named
volume}. Multi-stage build memisahkan tahap kompilasi dari tahap runtime sehingga
image akhir tidak membawa compiler atau \texttt{node\_modules}. Named volume
menyimpan data MySQL di luar siklus hidup container, sehingga data tidak hilang
saat container dihapus.

\section{FrankenPHP}

FrankenPHP adalah server aplikasi PHP yang dibangun di atas Caddy. Ia dapat
menjalankan PHP sebagai worker dan menangani HTTP secara langsung. Pada proyek ini,
FrankenPHP memakai konfigurasi Caddy untuk melayani direktori \texttt{public/}
Laravel pada port 8000. Pilihan ini berbeda dari \texttt{php artisan serve} yang
hanya untuk pengembangan, dan berbeda dari PHP-FPM yang memerlukan web server
terpisah.

FrankenPHP juga menulis access log ke stdout. Log tersebut dapat dibaca dengan
\texttt{docker compose logs backend} dan menjadi bukti bahwa permintaan benar-benar
melewati server tersebut.

\section{GitHub Actions, GHCR, dan Branch Protection}

GitHub Actions menjalankan workflow berdasarkan event. Pada proyek ini event yang
dipakai adalah \texttt{push} ke \texttt{main}, \texttt{dev}, dan \texttt{feature/**},
serta \texttt{pull\_request} ke \texttt{main} dan \texttt{dev}. Job
\texttt{backend-test} menjalankan \texttt{php artisan test}; job
\texttt{frontend-build} menjalankan \texttt{npm ci} dan \texttt{npm run build};
job \texttt{deploy} hanya berjalan pada push ke \texttt{main}.

GHCR (GitHub Container Registry) menyimpan image container. Image diberi tag commit
SHA, bukan \texttt{latest}, agar versi yang di-deploy dapat dilacak kembali ke
commit. Branch protection mencegah push langsung ke \texttt{main} dan \texttt{dev};
perubahan harus melalui pull request dan status check harus lulus.

\section{Arsitektur Sistem}

Gambar~\ref{fig:arsitektur} menunjukkan hubungan antar-komponen. Browser memuat
frontend dari Nginx pada port 8080. Permintaan ke \texttt{/api/} diteruskan Nginx ke
service \texttt{backend:8000}. Backend memakai FrankenPHP, memproses route Laravel,
dan mengakses MySQL melalui service \texttt{db:3306}. Semua service berada pada
network \texttt{app\_net}; nama service dipakai sebagai alamat, bukan IP container.

\begin{figure}[H]
\centering
\begin{tikzpicture}[
  node distance=1.1cm and 1.6cm,
  box/.style={draw, rounded corners, align=center, minimum width=3.1cm,
              minimum height=1.25cm, fill=blue!5, font=\small},
  data/.style={draw, cylinder, shape border rotate=90, aspect=0.25,
               align=center, minimum width=2.4cm, minimum height=1.4cm,
               fill=orange!10, font=\small},
  arrow/.style={-{Latex[length=2mm]}, thick}
]
\node[box] (browser) {Browser\\Vue 3 + Axios};
\node[box, right=of browser] (nginx) {Nginx\\frontend:80};
\node[box, right=of nginx] (backend) {FrankenPHP\\backend:8000};
\node[data, below=of backend] (mysql) {MySQL 8.4\\db:3306};
\node[box, below=of nginx] (volume) {Named volume\\\texttt{uts-pl\_db\_data}};

\draw[arrow] (browser) -- node[above,font=\scriptsize]{HTTP :8080} (nginx);
\draw[arrow] (nginx) -- node[above,font=\scriptsize]{/api/ proxy} (backend);
\draw[arrow] (backend) -- node[right,font=\scriptsize]{Eloquent} (mysql);
\draw[arrow, dashed] (mysql) -- (volume);
\node[font=\scriptsize, below=0.5cm of browser] {network \texttt{app\_net}};
\end{tikzpicture}
\caption{Arsitektur runtime aplikasi Buku Servis}
\label{fig:arsitektur}
\end{figure}

\section{Alur Data}

Alur data dari klik pengguna sampai data tersimpan dijelaskan pada
Gambar~\ref{fig:alur}. Contoh yang dipakai adalah penambahan catatan servis.

\begin{figure}[H]
\centering
\begin{tikzpicture}[
  node distance=0.5cm,
  proses/.style={draw, rounded corners, align=center, text width=2.45cm,
                 minimum height=1.35cm, font=\scriptsize, fill=gray!5},
  arrow/.style={-{Latex[length=2mm]}, thick}
]
\node[proses] (form) {Form Vue\\\texttt{ServiceForm}};
\node[proses, right=of form] (axios) {Axios\\\texttt{POST /api/}\allowbreak\texttt{service-logs}};
\node[proses, right=of axios] (nginx) {Nginx\\proxy \texttt{/api/}};
\node[proses, right=of nginx] (route) {Laravel route\\+ Sanctum};
\node[proses, right=of route] (valid) {FormRequest\\validasi};
\node[proses, below=of valid] (ctrl) {Controller\\\texttt{store()}};
\node[proses, left=of ctrl] (model) {Model\\\texttt{ServiceLog}};
\node[proses, left=of model] (sql) {MySQL\\tabel \texttt{service\_logs}};
\draw[arrow] (form) -- (axios);
\draw[arrow] (axios) -- (nginx);
\draw[arrow] (nginx) -- (route);
\draw[arrow] (route) -- (valid);
\draw[arrow] (valid) -- node[right,font=\tiny]{lolos} (ctrl);
\draw[arrow] (ctrl) -- node[above,font=\tiny]{\texttt{create()}} (model);
\draw[arrow] (model) -- node[above,font=\tiny]{INSERT} (sql);
\draw[arrow, dashed] (sql) -- node[below,font=\tiny]{201 + JSON} (form);
\end{tikzpicture}
\caption{Alur penambahan catatan servis dari Vue sampai MySQL}
\label{fig:alur}
\end{figure}

Jika validasi gagal, \texttt{FormRequest} menghentikan proses dan Laravel
mengembalikan 422 berisi \texttt{errors} per field. Jika token tidak ada, middleware
Sanctum mengembalikan 401 sebelum controller dijalankan. Jika token valid tetapi
catatan milik pengguna lain, query \texttt{findOrFail} menghasilkan 404. Pola ini
membuat respons dapat dibedakan oleh frontend dan oleh penguji API.

\chapter{ANALISIS DAN PERANCANGAN}
\label{chap:rancang}

\section{Kebutuhan Fungsional}

\begin{table}[H]
\centering
\caption{Kebutuhan fungsional aplikasi}
\label{tab:fungsi}
\small
\begin{tabularx}{\textwidth}{@{}p{0.8cm} X >{\raggedright\arraybackslash}p{4.6cm}@{}}
\toprule
\textbf{No} & \textbf{Kebutuhan} & \textbf{Bukti} \\
\midrule
F-01 & Pengguna dapat masuk dengan email dan kata sandi, lalu menerima token. & \texttt{POST}\newline\texttt{/login} \\
F-02 & Pengguna dapat keluar dan token aktif dicabut. & \texttt{POST}\newline\texttt{/logout} \\
F-03 & Pengguna dapat melihat daftar catatan servis miliknya. & \texttt{GET}\newline\texttt{/service-logs} \\
F-04 & Pengguna dapat menambah catatan servis. & \texttt{POST}\newline\texttt{/service-logs} \\
F-05 & Pengguna dapat melihat detail satu catatan. & \texttt{GET}\newline\texttt{/service-logs/\{id\}} \\
F-06 & Pengguna dapat mengubah catatan. & \texttt{PUT}\newline\texttt{/service-logs/\{id\}} \\
F-07 & Pengguna dapat menghapus catatan. & \texttt{DELETE}\newline\texttt{/service-logs/\{id\}} \\
F-08 & Data pengguna lain tidak dapat dibaca atau diubah. & query terbatas \texttt{user\_id} \\
F-09 & Frontend menampilkan ringkasan odometer dan total biaya. & halaman daftar \\
F-10 & Aplikasi menolak input tidak valid dengan pesan per field. & 422 + \texttt{errors} \\
\bottomrule
\end{tabularx}
\end{table}

Seluruh path pada kolom Bukti diawali \texttt{/api}.

\section{Kebutuhan Nonfungsional}

\begin{table}[H]
\centering
\caption{Kebutuhan nonfungsional dan cara pemenuhannya}
\label{tab:nonfungsi}
\begin{tabularx}{\textwidth}{@{}p{3.2cm} X@{}}
\toprule
\textbf{Aspek} & \textbf{Ketentuan} \\
\midrule
Keamanan & Token Sanctum; tidak ada registrasi publik; rahasia hanya di GitHub Secrets dan berkas \texttt{.env} yang tidak di-commit. \\
Keterlacakan & Image GHCR diberi tag commit SHA, bukan \texttt{latest}. \\
Portabilitas & Seluruh sistem berjalan dengan Docker Compose; tidak perlu PHP/Node di host. \\
Konsistensi data & Validasi server-side; tipe kolom MySQL; \texttt{user\_id} wajib. \\
Ketersediaan & Healthcheck backend \texttt{/up} dan frontend \texttt{/}; restart \texttt{unless-stopped}. \\
Kemudahan operasi & Satu perintah \texttt{docker compose up -d --build}; migrasi otomatis dari entrypoint. \\
Kualitas kode & PSR-4, FormRequest, service layer tipis, test fitur. \\
Dokumentasi & README, Postman Collection, komentar pada konfigurasi non-obvious. \\
\bottomrule
\end{tabularx}
\end{table}

\section{Model Data}

Aplikasi memakai tiga tabel utama: \texttt{users}, \texttt{service\_logs}, dan
\texttt{personal\_access\_tokens}. Tabel \texttt{users} dimiliki Laravel sejak awal.
Tabel \texttt{service\_logs} menyimpan catatan servis dan memiliki foreign key
\texttt{user\_id}. Tabel \texttt{personal\_access\_tokens} menyimpan token Sanctum.

\begin{table}[H]
\centering
\caption{Struktur tabel \texttt{service\_logs}}
\label{tab:servicelogs}
\begin{tabularx}{\textwidth}{@{}l l X@{}}
\toprule
\textbf{Kolom} & \textbf{Tipe} & \textbf{Keterangan} \\
\midrule
\texttt{id} & bigint unsigned & Primary key, auto increment. \\
\texttt{user\_id} & bigint unsigned & Foreign key ke \texttt{users.id}, cascade delete. \\
\texttt{motor} & varchar(100) & Nama atau tipe motor. \\
\texttt{serviced\_at} & date & Tanggal servis, format \texttt{YYYY-MM-DD}. \\
\texttt{odometer\_km} & int unsigned & Kilometer saat servis. \\
\texttt{service\_type} & varchar(100) & Jenis pekerjaan servis. \\
\texttt{cost} & int unsigned & Biaya dalam rupiah. \\
\texttt{workshop} & varchar(100) nullable & Nama bengkel. \\
\texttt{notes} & text nullable & Catatan tambahan. \\
\texttt{created\_at}, \texttt{updated\_at} & timestamp & Audit waktu. \\
\bottomrule
\end{tabularx}
\end{table}

Relasi yang dipakai: satu \texttt{User} memiliki banyak \texttt{ServiceLog};
satu \texttt{ServiceLog} dimiliki satu \texttt{User}. Indeks gabungan
\texttt{(user\_id, serviced\_at)} membantu query daftar yang selalu difilter
berdasarkan pengguna dan diurutkan berdasarkan tanggal.

\section{Rancangan API}

API mengikuti gaya REST dengan prefix \texttt{/api}. Tabel~\ref{tab:api}
merangkum endpoint, autentikasi, dan status yang diharapkan.

\begin{table}[H]
\centering
\caption{Rancangan endpoint API}
\label{tab:api}
\small
\begin{tabularx}{\textwidth}{@{}l l X l@{}}
\toprule
\textbf{Method} & \textbf{Endpoint} & \textbf{Keterangan} & \textbf{Status} \\
\midrule
POST & \texttt{/api/login} & Autentikasi, menghasilkan token & 200/401/422 \\
GET & \texttt{/api/me} & Profil pemilik token & 200/401 \\
POST & \texttt{/api/logout} & Mencabut token aktif & 200/401 \\
GET & \texttt{/api/service-logs} & Daftar catatan milik pengguna & 200/401 \\
POST & \texttt{/api/service-logs} & Menambah catatan & 201/401/422 \\
GET & \texttt{/api/service-logs/\{id\}} & Detail catatan milik pengguna & 200/401/404 \\
PUT & \texttt{/api/service-logs/\{id\}} & Mengubah catatan & 200/401/404/422 \\
DELETE & \texttt{/api/service-logs/\{id\}} & Menghapus catatan & 200/401/404 \\
\bottomrule
\end{tabularx}
\end{table}

Aturan validasi yang dipakai:

\begin{itemize}[leftmargin=1.2cm]
  \item \texttt{motor}: wajib, teks, maksimal 100 karakter.
  \item \texttt{serviced\_at}: wajib, format tanggal \texttt{Y-m-d}.
  \item \texttt{odometer\_km}: wajib, bilangan bulat, 0--9.999.999.
  \item \texttt{service\_type}: wajib, teks, maksimal 100 karakter.
  \item \texttt{cost}: wajib, bilangan bulat, 0--999.999.999.
  \item \texttt{workshop}: opsional, teks, maksimal 100 karakter.
  \item \texttt{notes}: opsional, teks, maksimal 1.000 karakter.
\end{itemize}

Pesan validasi disimpan di \texttt{App\textbackslash Support\textbackslash ValidationMessages}
agar controller dan FormRequest memakai pesan yang sama. Untuk aturan ukuran
(\texttt{min}/\texttt{max}) pada kolom numerik, pesan diberikan per atribut karena
Laravel mengambil pesan size dari translator, bukan dari larik \texttt{messages()}
biasa. Detail ini sempat menjadi kendala dan dibahas pada Subbab~\ref{sec:kendala}.

\section{Rancangan Antarmuka}

Ada dua halaman utama:

\begin{enumerate}[leftmargin=1.2cm]
  \item \textbf{Halaman masuk} (\texttt{/masuk}) berisi email, kata sandi, pesan
        error 401/422, dan tombol masuk. Halaman ini tidak menampilkan tautan
        registrasi karena registrasi publik ditutup.
  \item \textbf{Halaman catatan} (\texttt{/}) berisi ringkasan odometer, total
        biaya, formulir catat/ubah, dan daftar riwayat servis. Setiap baris memiliki
        tombol Ubah dan Hapus. Hapus memakai konfirmasi dua langkah.
\end{enumerate}

Komponen \texttt{ServiceForm} dipakai untuk tambah dan ubah. Saat prop
\texttt{editing} berisi catatan, formulir terisi data tersebut dan tombol berubah
menjadi ``Simpan perubahan''. Saat batal, formulir kembali kosong. Ringkasan
odometer memakai komponen \texttt{Odometer} yang menampilkan angka besar dan
total biaya dihitung dari daftar yang sedang tampil.

\section{Rancangan Deployment}

\begin{table}[H]
\centering
\caption{Service pada Docker Compose}
\label{tab:compose}
\begin{tabularx}{\textwidth}{@{}l l X@{}}
\toprule
\textbf{Service} & \textbf{Image} & \textbf{Peran} \\
\midrule
\texttt{db} & \texttt{mysql:8.4.13} & Basis data, volume \texttt{uts-pl\_db\_data}, healthcheck \texttt{mysqladmin ping}. \\
\texttt{backend} & \texttt{uts-pl-backend:local} & Laravel + FrankenPHP pada port 8000, healthcheck \texttt{/up}. \\
\texttt{frontend} & \texttt{uts-pl-frontend:local} & Nginx pada port 80, diteruskan ke host 8080, proxy \texttt{/api/}. \\
\bottomrule
\end{tabularx}
\end{table}

Kredensial MySQL diambil dari \texttt{.env} yang tidak di-commit. Berkas
\texttt{.env.example} menyediakan nama variabel tanpa nilai rahasia. Backend
menunggu \texttt{db} sehat sebelum mulai, lalu entrypoint menjalankan migrasi
dengan percobaan ulang. Frontend dibangun dengan argumen \texttt{VITE\_API\_URL=/api},
sehingga browser memanggil origin yang sama dan tidak ada masalah CORS.

\section{Rancangan CI/CD dan Strategi Branch}

\begin{table}[H]
\centering
\caption{Perilaku pipeline berdasarkan event}
\label{tab:pipeline}
\begin{tabularx}{\textwidth}{@{}l l X@{}}
\toprule
\textbf{Event} & \textbf{Job yang berjalan} & \textbf{Deploy} \\
\midrule
Push ke \texttt{feature/**} & \texttt{backend-test}, \texttt{frontend-build} & Tidak \\
Pull request ke \texttt{dev} & \texttt{backend-test}, \texttt{frontend-build} & Tidak \\
Push ke \texttt{dev} & \texttt{backend-test}, \texttt{frontend-build} & Tidak \\
Pull request ke \texttt{main} & \texttt{backend-test}, \texttt{frontend-build} & Tidak \\
Push ke \texttt{main} & \texttt{backend-test}, \texttt{frontend-build}, \texttt{deploy} & Ya, GHCR \\
\bottomrule
\end{tabularx}
\end{table}

Branch yang dipakai: \texttt{main} sebagai rilis, \texttt{dev} sebagai integrasi,
dan \texttt{feature/docker-frankenphp} sebagai tempat pekerjaan fitur. Perubahan
masuk melalui pull request. Branch protection pada \texttt{main} dan \texttt{dev}
mewajibkan pull request, status check \texttt{backend-test} dan
\texttt{frontend-build}, branch up-to-date, serta berlaku untuk administrator.

\section{Skenario Pengujian}

Pengujian dilakukan pada lima lapisan:

\begin{enumerate}[leftmargin=1.2cm]
  \item \textbf{Unit/feature test} dengan \texttt{php artisan test} untuk autentikasi,
        otorisasi, validasi, dan CRUD.
  \item \textbf{API manual} dengan curl dan Postman untuk memeriksa status HTTP
        serta bentuk JSON.
  \item \textbf{Antarmuka} dengan browser: login, tambah, ubah, hapus, dan kembali
        ke halaman masuk setelah logout.
  \item \textbf{Container} dengan \texttt{docker compose ps}, log backend, dan uji
        persistensi volume.
  \item \textbf{Pipeline} dengan push ke branch, pull request, dan push ke
        \texttt{main}; status job dan image GHCR diperiksa.
\end{enumerate}

\chapter{IMPLEMENTASI}
\label{chap:implementasi}

\section{Lingkungan Pengembangan}

Pekerjaan dilakukan pada Windows 11 dengan Docker Desktop. Versi perkakas yang
tercatat saat pengerjaan ditunjukkan pada Tabel~\ref{tab:versi}. Versi di host
berbeda dari versi di container; yang menjadi acuan runtime adalah versi di dalam
image.

\begin{table}[H]
\centering
\caption{Versi perkakas pengembangan dan runtime}
\label{tab:versi}
\begin{tabularx}{\textwidth}{@{}l l X@{}}
\toprule
\textbf{Perkakas} & \textbf{Versi} & \textbf{Keterangan} \\
\midrule
Docker Engine & 28.5.1 & Docker Desktop di Windows. \\
Docker Compose & v2.40.3-desktop.1 & Perintah \texttt{docker compose}. \\
Git & 2.45.1 & Windows. \\
PHP host & 8.2.12 & Untuk menjalankan tes lokal. \\
Composer host & 2.8.11 & Untuk dependensi lokal. \\
Node host & 22.19.0 & Untuk build frontend lokal. \\
PHP container & 8.4.27 & Dari image FrankenPHP. \\
Composer container & 2.10.3 & Dari image \texttt{composer:2.10.3}. \\
Node container & 22.23.4 & Tahap build frontend. \\
Nginx container & 1.30.6 & Tahap runtime frontend. \\
MySQL & 8.4.13 & Service \texttt{db}. \\
\bottomrule
\end{tabularx}
\end{table}

\section{Struktur Repositori}

Repositori dipisahkan menjadi \texttt{backend}, \texttt{frontend}, \texttt{postman},
\texttt{laporan}, dan berkas orkestrasi di akar. Pemisahan ini membuat build context
Docker tetap kecil dan mencegah frontend ikut masuk ke image backend.

\begin{lstlisting}[language=bash,caption={Struktur direktori utama repositori},label={kode:struktur}]
.
|-- .github/workflows/ci.yml
|-- backend/                 # Laravel 12, Sanctum, FrankenPHP
|   |-- app/Http/Controllers
|   |-- app/Http/Requests
|   |-- app/Models
|   |-- docker/Caddyfile
|   |-- docker/entrypoint.sh
|   |-- Dockerfile
|   |-- routes/api.php
|   `-- tests/Feature
|-- frontend/                # Vue 3, Vite, Axios, Nginx
|   |-- docker/nginx.conf
|   |-- src/api.js
|   |-- src/views
|   |-- Dockerfile
|   `-- package.json
|-- postman/buku-servis.postman_collection.json
|-- laporan/
|-- docker-compose.yml
|-- .env.example
`-- README.md
\end{lstlisting}

\section{Backend Laravel}

\subsection{Routing dan Perlindungan Token}

Berkas \texttt{routes/api.php} sengaja pendek. Hanya login yang berada di luar grup
\texttt{auth:sanctum}. Semua route lain, termasuk \texttt{/me} dan
\texttt{/logout}, berada di dalam grup. Tidak ada route registrasi.

\begin{lstlisting}[language=PHP,caption={Routing API di \texttt{backend/routes/api.php}},label={kode:routes}]
<?php

use App\Http\Controllers\AuthController;
use App\Http\Controllers\ServiceLogController;
use Illuminate\Support\Facades\Route;

// Hanya login yang dapat diakses tanpa token.
Route::post('/login', [AuthController::class, 'login']);

Route::middleware('auth:sanctum')->group(function () {
    Route::get('/me', [AuthController::class, 'me']);
    Route::post('/logout', [AuthController::class, 'logout']);

    Route::apiResource('service-logs', ServiceLogController::class);
});
\end{lstlisting}

\subsection{Autentikasi}

\texttt{AuthController::login()} memvalidasi email dan kata sandi, mencari pengguna,
dan memeriksa hash kata sandi. Jika gagal, respons 401 dengan pesan yang tidak
membocorkan apakah email atau kata sandi yang salah. Jika berhasil, token dibuat
dengan nama \texttt{spa} dan dikirim bersama data pengguna.

\begin{lstlisting}[language=PHP,caption={Login dan logout pada \texttt{AuthController}},label={kode:auth}]
public function login(Request $request): JsonResponse
{
    $data = $request->validate([
        'email' => ['required', 'email'],
        'password' => ['required', 'string'],
    ], ValidationMessages::messages(), ValidationMessages::attributes());

    $user = User::where('email', $data['email'])->first();

    if (! $user || ! Hash::check($data['password'], $user->password)) {
        return response()->json(['message' => 'Email atau kata sandi salah.'], 401);
    }

    return response()->json([
        'user' => $user,
        'token' => $user->createToken('spa')->plainTextToken,
    ]);
}

public function logout(Request $request): JsonResponse
{
    $request->user()->currentAccessToken()->delete();

    return response()->json(['message' => 'Berhasil keluar.']);
}
\end{lstlisting}

\subsection{CRUD dan Kepemilikan Data}

\texttt{ServiceLogController} memakai relasi \texttt{\$request->user()->serviceLogs()}
untuk semua query. Cara ini lebih aman daripada menerima \texttt{user\_id} dari
request karena pengguna tidak dapat memilih pemilik catatan. Method
\texttt{findOwned()} memakai \texttt{findOrFail} pada relasi, sehingga catatan milik
pengguna lain menghasilkan 404, bukan 403 yang membocorkan keberadaan data.

\begin{lstlisting}[language=PHP,caption={Inti CRUD pada \texttt{ServiceLogController}},label={kode:crud}]
public function index(Request $request): JsonResponse
{
    $logs = $request->user()
        ->serviceLogs()
        ->orderByDesc('serviced_at')
        ->orderByDesc('id')
        ->get();

    return response()->json(['data' => $logs]);
}

public function store(ServiceLogRequest $request): JsonResponse
{
    $log = $request->user()->serviceLogs()->create($request->validated());

    return response()->json(['data' => $log], 201);
}

public function update(ServiceLogRequest $request, string $id): JsonResponse
{
    $log = $this->findOwned($request, $id);
    $log->update($request->validated());

    return response()->json(['data' => $log]);
}

public function destroy(Request $request, string $id): JsonResponse
{
    $this->findOwned($request, $id)->delete();

    return response()->json(['message' => 'Catatan dihapus.']);
}

private function findOwned(Request $request, string $id)
{
    return $request->user()->serviceLogs()->findOrFail($id);
}
\end{lstlisting}

\subsection{Validasi}

Validasi ditempatkan di \texttt{ServiceLogRequest} agar controller tetap ringkas
dan aturan yang sama dipakai untuk \texttt{store} dan \texttt{update}. Pesan
berbahasa Indonesia dipusatkan di \texttt{ValidationMessages}.

\begin{lstlisting}[language=PHP,caption={Aturan validasi catatan servis},label={kode:validasi}]
public function rules(): array
{
    return [
        'motor' => ['required', 'string', 'max:100'],
        'serviced_at' => ['required', 'date_format:Y-m-d'],
        'odometer_km' => ['required', 'integer', 'min:0', 'max:9999999'],
        'service_type' => ['required', 'string', 'max:100'],
        'cost' => ['required', 'integer', 'min:0', 'max:999999999'],
        'workshop' => ['nullable', 'string', 'max:100'],
        'notes' => ['nullable', 'string', 'max:1000'],
    ];
}
\end{lstlisting}

\subsection{Migrasi dan Seeder}

Migrasi membuat tabel \texttt{service\_logs} dengan foreign key dan indeks gabungan.
\texttt{DemoSeeder} membuat akun \demoakun{} dan lima catatan contoh. Seeder
dipanggil dari entrypoint, bukan dari route registrasi. Dengan begitu, aplikasi
tetap memenuhi aturan ``semua endpoint selain login memerlukan token'' tanpa
mengorbankan kemudahan pengujian.

\begin{lstlisting}[language=PHP,caption={Potongan migrasi \texttt{service\_logs}},label={kode:migrasi}]
Schema::create('service_logs', function (Blueprint $table) {
    $table->id();
    $table->foreignId('user_id')->constrained()->cascadeOnDelete();
    $table->string('motor', 100);
    $table->date('serviced_at');
    $table->unsignedInteger('odometer_km');
    $table->string('service_type', 100);
    $table->unsignedInteger('cost');
    $table->string('workshop', 100)->nullable();
    $table->text('notes')->nullable();
    $table->timestamps();

    $table->index(['user_id', 'serviced_at']);
});
\end{lstlisting}

\subsection{Penanganan Permintaan Tanpa Token}

Saat pengujian container, ditemukan bahwa permintaan tanpa header
\texttt{Accept: application/json} dari Nginx sempat menghasilkan 500, bukan 401.
Penyebabnya adalah middleware \texttt{Authenticate} mencoba mengarahkan tamu ke
route bernama \texttt{login} yang tidak ada. Perbaikannya adalah memberi tahu
middleware bahwa aplikasi ini API-only dan tidak memiliki halaman login web.

\begin{lstlisting}[language=PHP,caption={Perbaikan redirect tamu pada \texttt{bootstrap/app.php}},label={kode:redirect}]
->withMiddleware(function (Middleware $middleware): void {
    // Aplikasi ini API-only. Tamu tidak boleh diarahkan ke route login web
    // yang tidak ada; biarkan exception handler mengembalikan 401 JSON.
    $middleware->redirectGuestsTo(fn () => null);
})
\end{lstlisting}

Perbaikan ini disertai tes regresi
\texttt{test\_guest\_without}\allowbreak\texttt{\_json\_accept}\allowbreak\texttt{\_header\_still\_gets\_401}.
Kasus ini
menunjukkan pentingnya menguji permintaan persis seperti yang dikirim proxy, bukan
hanya \texttt{getJson()} yang selalu menambahkan header JSON.

\section{Frontend Vue}

\subsection{Konfigurasi Alamat API}

Alamat API tidak ditulis di komponen. Nilainya berasal dari
\texttt{VITE\_API\_URL}. Di lokal, \texttt{frontend/.env} dapat berisi
\texttt{http://localhost:8000/api}; di Docker, build arg diisi \texttt{/api} dan
Nginx meneruskan permintaan ke \texttt{backend:8000}.

\begin{lstlisting}[language=JavaScript,caption={Instance Axios dan interceptor token},label={kode:axios}]
const api = axios.create({
  baseURL: import.meta.env.VITE_API_URL,
  headers: { Accept: 'application/json' },
})

api.interceptors.request.use((config) => {
  const token = localStorage.getItem('token')
  if (token) config.headers.Authorization = `Bearer ${token}`
  return config
})

api.interceptors.response.use(
  (response) => response,
  (error) => {
    const url = error.config?.url ?? ''
    const isAuthCall = url === '/login' || url === '/register'
    if (error.response?.status === 401 && !isAuthCall) {
      window.dispatchEvent(new Event('auth:expired'))
    }
    return Promise.reject(error)
  },
)
\end{lstlisting}

\subsection{State Autentikasi dan Router}

Token dan data pengguna disimpan di \texttt{localStorage} agar sesi bertahan setelah
halaman dimuat ulang. \texttt{auth.login()} memanggil API, menyimpan token, lalu
router mengarahkan ke halaman catatan. \texttt{auth.logout()} tetap menghapus token
lokal walaupun permintaan logout ke server gagal.

\begin{lstlisting}[language=JavaScript,caption={Guard router dan penyimpanan sesi},label={kode:router}]
router.beforeEach((to) => {
  if (to.meta.auth && !auth.isLoggedIn) return '/masuk'
  if (to.meta.guest && auth.isLoggedIn) return '/'
})

function save({ token, user }) {
  state.token = token
  state.user = user
  localStorage.setItem('token', token)
  localStorage.setItem('user', JSON.stringify(user))
}
\end{lstlisting}

\subsection{Formulir dan Daftar Catatan}

Formulir memakai \texttt{watch} pada prop \texttt{editing}. Saat pengguna menekan
Ubah, data catatan disalin ke form. Saat simpan, payload dikirim dengan method
\texttt{PUT} jika sedang mengubah, atau \texttt{POST} jika menambah. Respons 422
dipetakan ke \texttt{errors} per field; respons 401 diserahkan ke interceptor.

Daftar catatan diurutkan dari tanggal terbaru. Ringkasan menampilkan odometer
tertinggi dan total biaya dari catatan yang tampil. Konfirmasi hapus tidak memakai
\texttt{window.confirm}, tetapi state \texttt{confirmingId} agar tampilan konsisten
dan dapat diuji secara visual.

\section{Docker}

\subsection{Backend: FrankenPHP Multi-stage}

Dockerfile backend memiliki dua stage. Stage \texttt{vendor} memasang dependensi
Composer. Stage \texttt{runtime} memakai image FrankenPHP yang sama, membuat user
non-root \texttt{app}, menyalin hasil build, lalu menjalankan entrypoint. Perintah
\texttt{composer install} sengaja dijalankan sebelum kode disalin agar layer
dependensi tetap ter-cache saat hanya kode yang berubah.

\begin{lstlisting}[language=bash,caption={Bagian penting Dockerfile backend},label={kode:backenddocker}]
FROM dunglas/frankenphp:1.13.3-php8.4-alpine AS vendor
RUN install-php-extensions pdo_mysql intl bcmath opcache
COPY --from=composer:2.10.3 /usr/bin/composer /usr/local/bin/composer
WORKDIR /app
COPY composer.json composer.lock ./
RUN composer install --no-dev --no-scripts --no-autoloader \
    --no-interaction --prefer-dist --no-progress
COPY . .
RUN composer dump-autoload --no-dev --optimize --classmap-authoritative

FROM dunglas/frankenphp:1.13.3-php8.4-alpine AS runtime
RUN install-php-extensions pdo_mysql intl bcmath opcache \
    && addgroup -S app && adduser -S -G app -h /app app
WORKDIR /app
COPY --from=vendor --chown=app:app /app /app
COPY --chown=app:app docker/Caddyfile /etc/caddy/Caddyfile
COPY --chown=app:app docker/entrypoint.sh /usr/local/bin/entrypoint
USER app
EXPOSE 8000
HEALTHCHECK CMD wget -q -O /dev/null http://127.0.0.1:8000/up || exit 1
ENTRYPOINT ["entrypoint"]
CMD ["frankenphp", "run", "--config", "/etc/caddy/Caddyfile"]
\end{lstlisting}

Konfigurasi Caddy dibuat minimal: root di \texttt{/app/public}, \texttt{php\_server},
kompresi, dan access log ke stdout. Tidak ada port admin yang dibuka.

\begin{lstlisting}[language=bash,caption={Konfigurasi FrankenPHP pada \texttt{docker/Caddyfile}},label={kode:caddy}]
{
    frankenphp
    order php_server before file_server
    admin off
}

:8000 {
    root * /app/public
    encode zstd br gzip
    log {
        output stdout
        format console
    }
    php_server
}
\end{lstlisting}

Entrypoint menunggu MySQL siap, menjalankan migrasi dengan maksimal 30 percobaan,
menjalankan \texttt{DemoSeeder}, lalu membuat cache konfigurasi dan route. Pola ini
membuat container backend tidak langsung gagal ketika MySQL belum menerima koneksi.

\begin{lstlisting}[language=bash,caption={Entrypoint backend},label={kode:entrypoint}]
#!/bin/sh
set -eu

attempt=1
until php artisan migrate --force --no-interaction; do
    if [ "$attempt" -ge 30 ]; then
        echo "Migrasi gagal setelah $attempt percobaan." >&2
        exit 1
    fi
    echo "Basis data belum siap (percobaan $attempt/30), menunggu 2 detik..." >&2
    attempt=$((attempt + 1))
    sleep 2
done

php artisan db:seed --class=DemoSeeder --force --no-interaction
php artisan config:cache
php artisan route:cache

exec "$@"
\end{lstlisting}

\subsection{Frontend: Node Build dan Nginx Runtime}

Frontend memakai dua stage. Stage \texttt{build} memakai Node 22.23.4, menyalin
\texttt{package.json} dan \texttt{package-lock.json} lebih dulu, menjalankan
\texttt{npm ci}, lalu \texttt{npm run build}. Stage runtime hanya menyalin
\texttt{dist} ke Nginx. Tidak ada \texttt{node\_modules} di image akhir.

\begin{lstlisting}[language=bash,caption={Dockerfile frontend dua tahap},label={kode:frontenddocker}]
FROM node:22.23.4-alpine AS build
WORKDIR /app
COPY package.json package-lock.json ./
RUN npm ci --no-audit --no-fund
COPY . .
ARG VITE_API_URL=/api
ENV VITE_API_URL=$VITE_API_URL
RUN npm run build

FROM nginx:1.30.6-alpine3.24 AS runtime
COPY docker/nginx.conf /etc/nginx/conf.d/default.conf
COPY --from=build /app/dist /usr/share/nginx/html
EXPOSE 80
HEALTHCHECK CMD wget -q -O /dev/null http://127.0.0.1/ || exit 1
\end{lstlisting}

Nginx memakai \texttt{resolver 127.0.0.11} dan variabel \texttt{\$backend} agar
alamat service \texttt{backend} diperbarui melalui DNS internal Docker. Penulisan
IP container di konfigurasi dihindari karena IP dapat berubah saat container
dibuat ulang.

\begin{lstlisting}[language=bash,caption={Proxy \texttt{/api/} pada Nginx},label={kode:nginx}]
resolver 127.0.0.11 valid=10s ipv6=off;
set $backend http://backend:8000;

location /api/ {
    proxy_pass $backend$request_uri;
    proxy_http_version 1.1;
    proxy_set_header Host $host;
    proxy_set_header X-Real-IP $remote_addr;
    proxy_set_header X-Forwarded-For $proxy_add_x_forwarded_for;
    proxy_set_header X-Forwarded-Proto $scheme;
}
\end{lstlisting}

\subsection{Docker Compose}

Berkas \texttt{docker-compose.yml} mendefinisikan tiga service dan satu volume.
Backend menunggu \texttt{db} berstatus sehat. Frontend menunggu backend, tetapi
tetap dapat mulai sebelum backend siap karena Nginx akan mengembalikan 502 sampai
backend hidup; healthcheck frontend hanya memeriksa halaman statis.

\begin{lstlisting}[language=bash,caption={Ringkasan service pada \texttt{docker-compose.yml}},label={kode:compose}]
name: uts-pl

services:
  db:
    image: mysql:8.4.13
    environment:
      MYSQL_DATABASE: ${DB_DATABASE}
      MYSQL_USER: ${DB_USERNAME}
      MYSQL_PASSWORD: ${DB_PASSWORD}
      MYSQL_ROOT_PASSWORD: ${DB_ROOT_PASSWORD}
    volumes:
      - db_data:/var/lib/mysql
    healthcheck:
      test: ["CMD-SHELL", "mysqladmin ping -h 127.0.0.1 -u root -p\"$$MYSQL_ROOT_PASSWORD\" --silent"]

  backend:
    build: ./backend
    depends_on:
      db:
        condition: service_healthy
    environment:
      DB_HOST: db
      DB_DATABASE: ${DB_DATABASE}
      DB_USERNAME: ${DB_USERNAME}
      DB_PASSWORD: ${DB_PASSWORD}
    ports:
      - "8000:8000"

  frontend:
    build:
      context: ./frontend
      args:
        VITE_API_URL: /api
    ports:
      - "8080:80"

volumes:
  db_data:
\end{lstlisting}

\subsection{\texttt{.dockerignore} dan Rahasia}

Setiap build context memiliki \texttt{.dockerignore}. Backend mengabaikan
\texttt{vendor}, \texttt{.env}, \texttt{storage/logs}, dan cache. Frontend
mengabaikan \texttt{node\_modules}, \texttt{dist}, dan \texttt{.env}. Tidak ada
\texttt{ARG} atau \texttt{ENV} berisi kata sandi di Dockerfile. Nilai rahasia masuk
saat runtime melalui environment Compose yang dibaca dari \texttt{.env} yang tidak
di-commit.

\section{CI/CD}

Workflow \texttt{.github/workflows/ci.yml} memiliki tiga job. Job
\texttt{backend-test} memakai PHP 8.4, ekstensi SQLite, dan Composer; job
\texttt{frontend-build} memakai Node 22 dan \texttt{npm ci}; job \texttt{deploy}
memakai Buildx, login ke GHCR, lalu mendorong dua image dengan tag
\texttt{\$\{\{ github.sha \}\}}.

\begin{lstlisting}[language=bash,caption={Kondisi job deploy},label={kode:deploy}]
deploy:
  name: deploy
  if: github.event_name == 'push' && github.ref == 'refs/heads/main'
  needs:
    - backend-test
    - frontend-build
  runs-on: ubuntu-latest
\end{lstlisting}

Nama image dihitung dari \texttt{GITHUB\_REPOSITORY} dengan huruf kecil. Karena
repository bernama \texttt{Ryzen0000001/uts-pl-24-535196-SV-24171}, nama image
menjadi \ghcrbackend{} dan \ghcrfrontend{}. Tag SHA membuat setiap deploy dapat
dihubungkan ke commit yang tepat.

\begin{lstlisting}[language=bash,caption={Penentuan nama image GHCR},label={kode:ghcr}]
echo "BACKEND_IMAGE=ghcr.io/${GITHUB_REPOSITORY,,}-backend" >> "$GITHUB_ENV"
echo "FRONTEND_IMAGE=ghcr.io/${GITHUB_REPOSITORY,,}-frontend" >> "$GITHUB_ENV"
\end{lstlisting}

Rahasia yang disimpan di GitHub Secrets: \texttt{APP\_KEY}, \texttt{DB\_DATABASE},
\texttt{DB\_USERNAME}, \texttt{DB\_PASSWORD}, dan \texttt{DB\_ROOT\_PASSWORD}.
Tidak ada nilai rahasia di dalam repository, Dockerfile, atau workflow.

\section{Cara Menjalankan}

Langkah reproduksi dari repositori:

\begin{lstlisting}[language=bash,caption={Menjalankan seluruh sistem dari nol},label={kode:jalankan}]
git clone https://github.com/Ryzen0000001/uts-pl-24-535196-SV-24171.git
cd uts-pl-24-535196-SV-24171
cp .env.example .env
# isi APP_KEY, DB_PASSWORD, dan DB_ROOT_PASSWORD
docker compose up -d --build
docker compose ps
\end{lstlisting}

Setelah ketiga container sehat, frontend dapat dibuka di
\texttt{http://localhost:8080} dan API di \texttt{http://localhost:8000/api}.
Akun demo: \demoakun.

\chapter{HASIL DAN PEMBAHASAN}
\label{chap:hasil}

\section{Metode Pengujian}

Pengujian tidak hanya memeriksa ``berhasil'' atau ``gagal''. Setiap bukti diambil
dari eksekusi nyata dan disimpan di direktori \texttt{laporan/}. Untuk API, status
HTTP dan potongan body dicatat. Untuk container, keluaran \texttt{docker compose ps}
dan log backend disimpan. Untuk pipeline, tautan run dan tangkapan layar halaman
Actions diambil setelah run selesai.

\section{CPMK 1: REST API, Autentikasi, dan CRUD}

\subsection{Hasil Tes Otomatis}

Perintah \texttt{php artisan test} dijalankan di folder \texttt{backend}. Hasil
terakhir: \textbf{16 tes lulus, 38 asersi}. Tes mencakup registrasi yang tidak
tersedia, login valid dan tidak valid, logout yang mencabut token, akses tamu ke
semua endpoint, kepemilikan data, validasi, serta CRUD.

\begin{lstlisting}[language=bash,caption={Keluaran \texttt{php artisan test} (ringkas)},label={kode:test}]
PASS  Tests\Unit\ExampleTest
PASS  Tests\Feature\AuthTest
  OK register endpoint is not exposed
  OK login with invalid email returns 422
  OK login with valid credentials returns token
  OK login with wrong password returns 401
  OK logout revokes the token
  OK me without token returns 401
PASS  Tests\Feature\ExampleTest
PASS  Tests\Feature\ServiceLogTest
  OK guest gets 401 on every endpoint
  OK guest without json accept header still gets 401
  OK index lists only own logs
  OK store creates log and returns 201
  OK store with invalid data returns 422
  OK show returns own log but 404 for other users log
  OK update changes the log
  OK destroy deletes the log

Tests: 16 passed (38 assertions)
Duration: 0.45s
\end{lstlisting}

Keluaran lengkap ada pada Lampiran~\ref{lam:test}. Tes
\texttt{guest without json accept header still gets 401} adalah tes yang ditambahkan
setelah menemukan bug 500 pada proxy Nginx. Tanpa tes ini, bug tersebut hanya akan
muncul saat diuji dari browser atau curl tanpa header JSON.

\subsection{Verifikasi API Manual}

Verifikasi dilakukan terhadap API yang berjalan di container, bukan server
pengembangan. Tabel~\ref{tab:hasilapi} merangkum hasilnya. Transkrip lengkap ada
pada Lampiran~\ref{lam:api}; token pada transkrip disamarkan.

\begin{table}[H]
\centering
\caption{Hasil verifikasi status HTTP API}
\label{tab:hasilapi}
\begin{tabularx}{\textwidth}{@{}l l X l@{}}
\toprule
\textbf{Method} & \textbf{Endpoint} & \textbf{Skenario} & \textbf{Status} \\
\midrule
POST & \texttt{/api/login} & Kredensial demo benar & 200 \\
POST & \texttt{/api/service-logs} & Payload valid & 201 \\
GET & \texttt{/api/service-logs} & Daftar milik pengguna & 200 \\
GET & \texttt{/api/service-logs/\{id\}} & Catatan milik sendiri & 200 \\
PUT & \texttt{/api/service-logs/\{id\}} & Ubah biaya & 200 \\
POST & \texttt{/api/service-logs} & \texttt{motor} kosong, \texttt{odometer} bukan angka, biaya negatif & 422 \\
GET & \texttt{/api/service-logs} & Tanpa header token & 401 \\
DELETE & \texttt{/api/service-logs/\{id\}} & Hapus catatan uji & 200 \\
POST & \texttt{/api/logout} & Cabut token & 200 \\
\bottomrule
\end{tabularx}
\end{table}

Contoh respons 422 dari pengujian:

\begin{lstlisting}[language=bash,caption={Contoh respons validasi 422},label={kode:422}]
{"message":"Motor wajib diisi. (and 3 more errors)",
 "errors":{
   "motor":["Motor wajib diisi."],
   "odometer_km":["Kilometer harus berupa angka bulat."],
   "service_type":["Jenis servis wajib diisi."],
   "cost":["Biaya tidak boleh kurang dari 0."]
 }}
\end{lstlisting}

\subsection{Postman Collection}

Berkas \texttt{postman/buku-servis.postman\_collection.json} berisi tiga folder:
Auth (Login, Me, Logout), Service Logs (List, Create, Create invalid, Show, Update,
Delete), dan Tanpa token (List tanpa token). Variabel koleksi \texttt{baseUrl}
dapat diubah ke \texttt{http://localhost:8000/api} atau
\texttt{http://localhost:8080/api}. Skrip tes pada Login menyimpan token ke variabel
koleksi; skrip pada Create menyimpan \texttt{logId}; skrip pada Create invalid dan
List tanpa token memeriksa status 422 dan 401.

\isipostman{postman-login.png}{Respons login di Postman dan penyimpanan token}
\isipostman{postman-create-201.png}{Respons pembuatan catatan dengan status 201}
\isipostman{postman-422.png}{Respons validasi dengan status 422}
\isipostman{postman-401.png}{Respons tanpa token dengan status 401}

\section{CPMK 2: Frontend Vue dan Integrasi Axios}

Antarmuka diuji dengan browser pada alamat \texttt{http://localhost:8080}. Alur yang
diuji: masuk dengan akun demo, melihat daftar, menambah catatan, mengubah catatan,
menghapus catatan, lalu keluar. Tangkapan layar diambil dari aplikasi yang berjalan
di container, bukan mockup.

\gambarpenting{img/login.png}{Halaman masuk sebelum autentikasi}
\label{fig:login}

\gambarpenting{img/dashboard.png}{Halaman daftar catatan setelah login}
\label{fig:dashboard}

Formulir tambah diisi dengan data uji: motor Vario 125, tanggal 24 Oktober 2026,
odometer 25.120 km, jenis servis ``Ganti oli gardan'', biaya Rp35.000, bengkel
AHASS Sleman, dan catatan ``Servis pagi, antre sebentar''. Sebelum disimpan,
ringkasan menunjukkan servis terakhir 12 Oktober 2026 dengan total Rp625.000 dari
lima catatan. Setelah disimpan, catatan muncul di bagian atas daftar karena
diurutkan berdasarkan tanggal terbaru; ringkasan berubah menjadi Rp660.000 dari
enam catatan.

\gambarpenting{img/tambah-form.png}{Formulir tambah catatan yang sudah diisi}
\label{fig:tambahform}

\gambarpenting{img/tambah-hasil.png}{Catatan baru muncul di daftar setelah disimpan}
\label{fig:tambahhasil}

Untuk mengubah, tombol ``Ubah'' pada baris catatan mengisi formulir dan mengubah
judul menjadi ``Ubah catatan''. Biaya diubah dari Rp35.000 menjadi Rp40.000 dan
catatan diganti menjadi ``Harga naik, tetap di AHASS''. Setelah disimpan, daftar
dimuat ulang dan total biaya pada ringkasan menjadi Rp665.000.

\gambarpenting{img/ubah-form.png}{Formulir ubah dengan data catatan terpilih}
\label{fig:ubahform}

\gambarpenting{img/ubah-hasil.png}{Daftar setelah biaya diubah}
\label{fig:ubahhasil}

Hapus memakai konfirmasi di dalam baris. Pengguna harus menekan ``Ya, hapus'';
tombol ``Batal'' menutup konfirmasi tanpa mengirim permintaan. Setelah catatan
Ganti oli gardan dihapus, daftar kembali berisi lima catatan dengan total
Rp625.000 dan pesan ``Catatan dihapus.'' muncul.

\gambarpenting{img/hapus-konfirmasi.png}{Konfirmasi hapus dua langkah}
\label{fig:hapuskonfirmasi}

\gambarpenting{img/hapus-hasil.png}{Daftar setelah catatan dihapus}
\label{fig:hapushasil}

Dari sisi integrasi, frontend memanggil \texttt{/api/...} pada origin yang sama
(\texttt{localhost:8080}). Nginx meneruskan ke backend. Karena token dikirim pada
header \texttt{Authorization}, tidak ada cookie lintas origin dan tidak ada
konfigurasi CORS khusus di frontend. Jika token dicabut, interceptor 401 memicu
kembali ke halaman masuk.

\section{CPMK 3: Docker Compose, FrankenPHP, dan Persistensi}

\subsection{Status Container}

Setelah \texttt{docker compose up -d --build}, ketiga container berjalan dan sehat.
Tabel~\ref{tab:status} merangkum status dari \texttt{docker compose ps}; keluaran
mentahnya ada pada Lampiran~\ref{lam:ps}.

\begin{table}[H]
\centering
\caption{Status container dari \texttt{docker compose ps}}
\label{tab:status}
\small
\begin{tabularx}{\textwidth}{@{}l l l X@{}}
\toprule
\textbf{Container} & \textbf{Service} & \textbf{Image} & \textbf{Status dan port} \\
\midrule
\texttt{uts-pl-backend-1} & backend & \texttt{uts-pl-backend:local} & Up (healthy), \texttt{0.0.0.0:8000->8000/tcp} \\
\texttt{uts-pl-db-1} & db & \texttt{mysql:8.4.13} & Up (healthy), port internal 3306 \\
\texttt{uts-pl-frontend-1} & frontend & \texttt{uts-pl-frontend:local} & Up (healthy), \texttt{0.0.0.0:8080->80/tcp} \\
\bottomrule
\end{tabularx}
\end{table}

\subsection{Bukti FrankenPHP}

Log backend memuat baris startup FrankenPHP dan versi PHP 8.4.27. Access log
menunjukkan permintaan \texttt{POST /api/login}, \texttt{POST /api/service-logs},
\texttt{GET /api/service-logs}, \texttt{PUT}, \texttt{DELETE}, dan respons 401
untuk permintaan tanpa token. Header respons juga memuat
\texttt{Server: FrankenPHP Caddy}. Ini membuktikan backend tidak dijalankan dengan
\texttt{php artisan serve} atau PHP-FPM.

\begin{lstlisting}[language=bash,caption={Potongan bukti log FrankenPHP},label={kode:franken}]
STARTUP: {"logger":"frankenphp","msg":"FrankenPHP started","php_version":"8.4.27",
          "total_threads":48,"worker_threads":0,"max_threads":48}
POST /api/login -> 200
POST /api/service-logs -> 201
GET /api/service-logs -> 200
GET /api/service-logs/7 -> 200
PUT /api/service-logs/7 -> 200
POST /api/service-logs -> 422
GET /api/service-logs -> 401
DELETE /api/service-logs/7 -> 200
POST /api/logout -> 200
\end{lstlisting}

Pemeriksaan proses di dalam container juga menunjukkan perintah
\texttt{frankenphp run --config /etc/caddy/Caddyfile} dan user \texttt{app}, bukan
root. Image backend berukuran sekitar 374 MB; image frontend sekitar 92,9 MB.
Image frontend tidak memuat \texttt{node\_modules}, dan image backend tidak memuat
Composer.

\subsection{Persistensi Volume}

Uji persistensi dilakukan dengan membuat catatan uji, menghentikan seluruh
container tanpa menghapus volume, lalu menyalakannya kembali. Catatan dengan id 8
masih ada setelah \texttt{docker compose up -d}. Setelah verifikasi, catatan uji
dihapus agar data demo tetap rapi.

\begin{lstlisting}[language=bash,caption={Hasil uji persistensi},label={kode:persistensi}]
SEBELUM DOWN: id uji = 8, jumlah catatan = 6
docker compose down
docker compose up -d
backend /up -> 200
SESUDAH UP: jumlah catatan = 6; data id=8 ditemukan = True;
              service_type = Uji persistensi volume
Pembersihan: DELETE id=8 -> {"message":"Catatan dihapus."}
\end{lstlisting}

\section{CPMK 4: Git, Branch Protection, dan CI/CD}

\subsection{Repositori dan Branch}

{\raggedright
Repositori publik berada di
\url{https://github.com/Ryzen0000001/uts-pl-24-535196-SV-24171}.
Branch yang ada: \texttt{main},
\texttt{dev}, dan \texttt{feature/docker-frankenphp}. Pekerjaan fitur masuk ke
\texttt{dev} melalui PR \#1, lalu \texttt{dev} masuk ke \texttt{main} melalui PR
\#2. Kedua PR sudah digabung dan branch fitur tetap disimpan sebagai jejak
pekerjaan.\par}

\gambarpenting[0.9\textwidth]{img/repo-home.png}{Halaman utama repositori publik}
\label{fig:repohome}

\gambarpenting[0.9\textwidth]{img/repo-branches.png}{Daftar branch \texttt{main}, \texttt{dev}, dan fitur}
\label{fig:repobranches}

\gambarpenting[0.9\textwidth]{img/repo-pulls.png}{Dua pull request yang sudah digabung}
\label{fig:repopulls}

\gambarpenting[0.82\textwidth]{img/pr-1.png}{PR \#1: fitur Docker dan FrankenPHP masuk ke \texttt{dev}}
\label{fig:pr1}

\gambarpenting[0.82\textwidth]{img/pr-2.png}{PR \#2: \texttt{dev} masuk ke \texttt{main}}
\label{fig:pr2}

\subsection{Branch Protection}

Aturan pada \texttt{main} dan \texttt{dev} mewajibkan pull request sebelum merge,
status check \texttt{backend-test} dan \texttt{frontend-build}, serta branch yang
up-to-date. Opsi ``Do not allow bypassing the above settings'' diaktifkan sehingga
aturan juga berlaku untuk administrator. Force push dan penghapusan branch
dinonaktifkan.

\gambarpenting[0.88\textwidth]{img/branch-protection-main-crop.png}{Branch protection pada \texttt{main}}
\label{fig:protmain}

\gambarpenting[0.88\textwidth]{img/branch-protection-dev-crop.png}{Branch protection pada \texttt{dev}}
\label{fig:protdev}

\subsection{GitHub Secrets}

Lima rahasia disimpan di \texttt{Settings > Secrets and variables > Actions}:
\texttt{APP\_KEY}, \texttt{DB\_DATABASE}, \texttt{DB\_USERNAME},
\texttt{DB\_PASSWORD}, dan \texttt{DB\_ROOT\_PASSWORD}. Halaman secrets hanya
menampilkan nama, bukan nilai.

\gambarpenting[0.78\textwidth]{img/github-secrets-crop.png}{Daftar GitHub Secrets tanpa nilai rahasia}
\label{fig:secrets}

\subsection{Hasil Pipeline}

Push ke branch fitur dan \texttt{dev} menjalankan dua job pengujian tanpa deploy.
Push ke \texttt{main} menjalankan ketiga job; job deploy membangun dan mendorong
image ke GHCR. Tabel~\ref{tab:runs} merangkum run yang diamati.

\begin{table}[H]
\centering
\caption{Run GitHub Actions yang diamati}
\label{tab:runs}
\small
\begin{tabularx}{\textwidth}{@{}l X l l@{}}
\toprule
\textbf{Event} & \textbf{Commit/Run} & \textbf{Job} & \textbf{Deploy} \\
\midrule
Push fitur & CI pada branch \texttt{feature/docker-frankenphp} & test + build & Tidak \\
PR \#1 & Pemeriksaan sebelum merge ke \texttt{dev} & test + build & Tidak \\
Push \texttt{dev} & Setelah PR \#1 digabung & test + build & Tidak \\
PR \#2 & Pemeriksaan sebelum merge ke \texttt{main} & test + build & Tidak \\
Push \texttt{main} & Setelah PR \#2 digabung & test + build + deploy & Ya \\
\bottomrule
\end{tabularx}
\end{table}

\gambarpenting[0.9\textwidth]{img/actions-dev-crop.png}{Run pada \texttt{dev}: dua job lulus, deploy tidak muncul}
\label{fig:actionsdev}

\gambarpenting[0.9\textwidth]{img/actions-main-crop.png}{Run pada \texttt{main}: backend-test, frontend-build, dan deploy lulus}
\label{fig:actionsmain}

\subsection{Image di GHCR}

Dua package container dipublikasikan: backend dan frontend. Halaman package
menampilkan tag commit SHA dan tautan ke repositori. Setelah repo berganti nama
menjadi \texttt{uts-pl}, nama package mengikuti nama repositori baru.

\gambarpenting[0.78\textwidth]{img/packages.png}{Daftar package container di GHCR}
\label{fig:packages}

\gambarpenting[0.8\textwidth]{img/package-backend-crop.png}{Halaman package backend dan tag SHA}
\label{fig:packagebackend}

\section{CPMK 5: Dokumentasi dan Reproduksibilitas}

Dokumentasi terdiri dari README, Postman Collection, komentar pada konfigurasi
non-obvious, dan laporan ini. README memuat cara menjalankan, tabel endpoint,
perintah pengujian, penjelasan alamat API, serta daftar secrets. Postman Collection
memungkinkan penguji mencoba endpoint tanpa menulis ulang body JSON.

Yang belum sepenuhnya otomatis adalah pengambilan tangkapan layar Postman dan
pengujian UI. Keduanya masih manual karena memerlukan aplikasi desktop dan
pemeriksaan visual. Untuk pengembangan lanjutan, UI test dapat memakai Playwright,
sedangkan API test dapat dijalankan dengan Newman di pipeline.

\section{Pembahasan Alur Data End-to-End}

Saat pengguna mengisi formulir dan menekan ``Simpan catatan'', Vue membentuk objek
payload. \texttt{ServiceForm} memanggil \texttt{api.post('/service-logs', payload)}.
Axios menambahkan \texttt{baseURL} dari \texttt{VITE\_API\_URL} dan header
\texttt{Authorization} dari token di \texttt{localStorage}. Di Docker, permintaan
menuju \texttt{http://localhost:8080/api/service-logs}. Nginx mencocokkan
\texttt{location /api/} dan meneruskannya ke \texttt{http://backend:8000/api/service-logs}.

Di backend, route Laravel mencocokkan \texttt{POST /api/service-logs}. Middleware
\texttt{auth:sanctum} memeriksa token. Jika valid, \texttt{ServiceLogRequest}
menjalankan validasi. Jika lolos, controller memanggil relasi
\texttt{\$request->user()->serviceLogs()->create(\$request->validated())}. Eloquent
mengisi \texttt{user\_id} dari pengguna yang terautentikasi dan menjalankan
\texttt{INSERT} ke MySQL. Controller mengembalikan JSON dengan status 201.

Frontend menerima respons, mengosongkan formulir, memuat ulang daftar dengan
\texttt{GET /api/service-logs}, dan menampilkan pesan ``Catatan disimpan.''. Jika
validasi gagal, respons 422 berisi \texttt{errors}; formulir menampilkan pesan di
bawah field yang salah. Jika token hilang, respons 401 memicu event
\texttt{auth:expired} dan router mengarahkan pengguna ke \texttt{/masuk}.

\section{Kendala dan Solusi}
\label{sec:kendala}

\begin{table}[H]
\centering
\caption{Kendala yang ditemui dan solusinya}
\label{tab:kendala}
\small
\begin{tabularx}{\textwidth}{@{}p{0.8cm} X X@{}}
\toprule
\textbf{No} & \textbf{Kendala} & \textbf{Solusi} \\
\midrule
1 & Frontend sempat memanggil \texttt{localhost:8000} dari browser, sehingga alamat API berbeda antara lokal dan Docker. & Alamat dipindah ke \texttt{VITE\_API\_URL}; Docker memakai \texttt{/api} dan Nginx mem-proxy ke \texttt{backend:8000}. \\
2 & Nginx dengan \texttt{proxy\_pass http://backend:8000} dapat menyimpan DNS lama saat container backend dibuat ulang. & Memakai \texttt{resolver 127.0.0.11} dan variabel \texttt{\$backend} agar resolusi dilakukan ulang. \\
3 & Backend sempat 500 untuk permintaan tanpa \texttt{Accept: application/json} karena middleware mencari route \texttt{login}. & Menambahkan \texttt{redirectGuestsTo(fn () => null)} dan tes regresi 401. \\
4 & Pesan validasi \texttt{min} untuk biaya masih berbahasa Inggris. & Menambahkan pesan per atribut \texttt{cost.min} dan \texttt{odometer\_km.min}; pesan size Laravel tidak dibaca dari larik biasa. \\
5 & Container backend mulai sebelum MySQL siap. & Entrypoint menunggu dengan 30 percobaan, lalu migrasi dan seeder. \\
6 & Registrasi publik tidak boleh ada, tetapi penguji membutuhkan akun. & \texttt{DemoSeeder} membuat akun demo saat backend mulai. \\
7 & Data hilang jika volume ikut dihapus. & Volume bernama \texttt{uts-pl\_db\_data}; dokumentasi menjelaskan perbedaan \texttt{down} dan \texttt{down -v}. \\
8 & Deploy bisa terjadi dari branch selain \texttt{main}. & Job deploy diberi kondisi \texttt{github.ref}\allowbreak\texttt{==}\allowbreak\texttt{'refs/heads/main'}; branch protection mencegah push langsung. \\
9 & Rahasia berisiko masuk repositori. & \texttt{.env} di-\texttt{.gitignore}; nilai produksi di GitHub Secrets; Dockerfile tidak memuat secret. \\
10 & Nama repositori awal memakai akhiran \texttt{uts-pi}, tidak sesuai kode mata kuliah. & Repositori, nama image, volume, README, dan laporan diseragamkan menjadi \texttt{uts-pl}. \\
\bottomrule
\end{tabularx}
\end{table}

Kendala nomor 3 adalah yang paling penting untuk dicatat. Tanpa pengujian dari
proxy, bug ini tidak terlihat karena \texttt{getJson()} selalu mengirim header
JSON. Perbaikannya kecil, tetapi dampaknya besar: API sekarang konsisten
mengembalikan 401 JSON untuk semua permintaan tanpa token.

\section{Kesesuaian dengan Ketentuan}

\begin{table}[H]
\centering
\caption{Pemetaan ketentuan tugas ke bukti}
\label{tab:checklist}
\small
\begin{tabularx}{\textwidth}{@{}p{4.3cm} X p{2.1cm}@{}}
\toprule
\textbf{Ketentuan} & \textbf{Bukti} & \textbf{Lokasi} \\
\midrule
REST API Laravel + Sanctum & \texttt{routes/api.php}, \texttt{AuthController}, tes 401 & Subbab 4.3, 5.2 \\
Semua endpoint selain login memakai token & Grup \texttt{auth:sanctum}; tes tamu 401 & Kode~\ref{kode:routes}, Tabel~\ref{tab:hasilapi} \\
CRUD status 201/422/401 & Tes dan transkrip curl & Tabel~\ref{tab:hasilapi} \\
Minimal 2 feature test & 16 tes, 38 asersi & Lampiran~\ref{lam:test} \\
Postman Collection & \texttt{postman/buku-servis.postman\_collection.json} & Subbab 5.2 \\
Vue + Axios + token & \texttt{api.js}, \texttt{auth.js}, halaman UI & Subbab 4.4, 5.3 \\
API URL dari env & \texttt{VITE\_API\_URL} & Kode~\ref{kode:axios} \\
Docker Compose 3 service & db, backend, frontend & Tabel~\ref{tab:compose} \\
Jaringan service-name & \texttt{DB\_HOST=db}, proxy \texttt{backend:8000} & Kode~\ref{kode:nginx} \\
Volume permanen & \texttt{uts-pl\_db\_data}; uji id 8 & Subbab 5.5 \\
Kredensial dari \texttt{.env} & \texttt{.env.example}, secrets & Subbab 4.7, 5.4 \\
FrankenPHP (bonus) & Dockerfile, Caddyfile, log startup & Subbab 4.5, 5.5 \\
CI/CD & \texttt{ci.yml}; job test, build, deploy & Tabel~\ref{tab:runs} \\
Deploy hanya dari \texttt{main} & Kondisi job deploy & Kode~\ref{kode:deploy} \\
Image GHCR tag SHA & Halaman package & Gambar~\ref{fig:packagebackend} \\
Branch + PR + protection & Branch, PR \#1, \#2, screenshot & Gambar~\ref{fig:protmain} \\
Secrets tidak di repo & GitHub Secrets, \texttt{.gitignore} & Gambar~\ref{fig:secrets} \\
\bottomrule
\end{tabularx}
\end{table}

\chapter{KESIMPULAN DAN SARAN}
\label{chap:kesimpulan}

\section{Kesimpulan}

Buku Servis berhasil dibangun sebagai aplikasi web terintegrasi. API Laravel 12
dengan Sanctum melindungi seluruh endpoint selain login; 16 tes fitur lulus dengan
38 asersi. Frontend Vue 3 memakai Axios dan mengambil alamat API dari
\texttt{VITE\_API\_URL}, sehingga kode yang sama dapat berjalan di lokal dan di
Docker. Docker Compose menjalankan MySQL, backend FrankenPHP, dan frontend Nginx
pada satu network. Data bertahan setelah container dihentikan karena memakai volume
bernama.

FrankenPHP terbukti dipakai: log startup mencatat PHP 8.4.27 dan access log mencatat
permintaan API, sedangkan header respons memuat \texttt{Server: FrankenPHP Caddy}.
Pipeline GitHub Actions menjalankan tes dan build pada branch fitur, \texttt{dev},
serta pull request; deploy hanya berjalan pada push ke \texttt{main} dan menghasilkan
image GHCR bertag commit SHA. Branch protection pada \texttt{main} dan \texttt{dev}
mewajibkan pull request serta status check, dan rahasia disimpan di GitHub Secrets.

Pekerjaan ini juga menunjukkan bahwa integrasi tidak selesai setelah kode CRUD
berjalan. Dua masalah nyata, yaitu DNS Nginx dan redirect tamu ke route login,
baru muncul saat sistem dijalankan sebagai container dan diuji dari proxy. Keduanya
diselesaikan dengan konfigurasi kecil dan tes regresi.

\section{Saran}

Untuk pengembangan berikutnya:

\begin{enumerate}[leftmargin=1.2cm]
  \item Tambahkan pengujian UI otomatis dengan Playwright agar alur login, tambah,
        ubah, dan hapus tidak hanya diuji manual.
  \item Jalankan Postman Collection dengan Newman di pipeline agar kontrak API
        diperiksa pada setiap pull request.
  \item Tambahkan fitur pengingat servis berdasarkan interval kilometer atau
        tanggal, serta ekspor catatan ke PDF/CSV.
  \item Pisahkan database production dari database demo dan gunakan migrasi
        terpisah untuk data contoh.
  \item Gunakan secret manager atau environment terenkripsi untuk deployment di
        server nyata, bukan hanya GitHub Secrets.
  \item Tambahkan rate limiting pada login dan audit log untuk permintaan yang
        mengubah data.
\end{enumerate}

% ---------- Daftar Pustaka ----------
\backmatter
\chapter*{DAFTAR PUSTAKA}
\addcontentsline{toc}{chapter}{DAFTAR PUSTAKA}
\phantomsection

\begin{thebibliography}{9}
\bibitem{laravel}
Laravel. \emph{Laravel 12.x Documentation}. \url{https://laravel.com/docs/12.x}, diakses 24 Oktober 2026.

\bibitem{sanctum}
Laravel. \emph{Laravel Sanctum}. \url{https://laravel.com/docs/12.x/sanctum}, diakses 24 Oktober 2026.

\bibitem{vue}
Vue.js. \emph{Vue 3 Documentation}. \url{https://vuejs.org/guide/}, diakses 24 Oktober 2026.

\bibitem{axios}
Axios. \emph{Axios Documentation}. \url{https://axios-http.com/docs/intro}, diakses 24 Oktober 2026.

\bibitem{docker}
Docker. \emph{Docker Compose Documentation}. \url{https://docs.docker.com/compose/}, diakses 24 Oktober 2026.

\bibitem{frankenphp}
FrankenPHP. \emph{FrankenPHP Documentation}. \url{https://frankenphp.dev/docs/}, diakses 24 Oktober 2026.

\bibitem{gha}
GitHub. \emph{GitHub Actions Documentation}. \url{https://docs.github.com/en/actions}, diakses 24 Oktober 2026.

\bibitem{ghcr}
GitHub. \emph{Working with the Container Registry}. \url{https://docs.github.com/en/packages/working-with-a-github-packages-registry/working-with-the-container-registry}, diakses 24 Oktober 2026.

\bibitem{protection}
GitHub. \emph{About protected branches}. \url{https://docs.github.com/en/repositories/configuring-branches-and-merges-in-your-repository/managing-protected-branches/about-protected-branches}, diakses 24 Oktober 2026.

\bibitem{postman}
Postman. \emph{Postman Collection Format}. \url{https://schema.getpostman.com/json/collection/v2.1.0/docs/index.html}, diakses 24 Oktober 2026.
\end{thebibliography}

% ---------- Lampiran ----------
\appendix
\clearpage
\phantomsection
\chapter*{LAMPIRAN}
\addcontentsline{toc}{chapter}{LAMPIRAN}
\markboth{LAMPIRAN}{}
% \backmatter membuat \chapter tidak menomori bab. Penomoran bagian
% lampiran diatur manual agar menjadi A.1, A.2, dan seterusnya.
\setcounter{chapter}{1}
\renewcommand{\thechapter}{A}
\setcounter{section}{0}

\section{Keluaran Tes}
\label{lam:test}
\lstinputlisting[language=bash,caption={Keluaran lengkap \texttt{php artisan test}}]{test-output.txt}

\section{Keluaran \texttt{docker compose ps}}
\label{lam:ps}
\lstinputlisting[language=bash,caption={Keluaran \texttt{docker compose ps} setelah semua container sehat}]{docker-compose-ps.txt}

\section{Bukti Log FrankenPHP}
\label{lam:franken}
\lstinputlisting[language=bash,caption={Startup dan access log FrankenPHP (tanpa warna)}]{frankenphp-evidence.txt}

\section{Bukti Persistensi}
\label{lam:persistensi}
\lstinputlisting[language=bash,caption={Transkrip uji persistensi volume}]{persistence.txt}

\section{Transkrip Verifikasi API}
\label{lam:api}
\lstinputlisting[language=bash,caption={Transkrip curl; token disamarkan},breaklines=true]{api-verification.txt}

\section{Tautan Penting}

\begin{itemize}[leftmargin=1.2cm]
  \item Repositori: \url{https://github.com/Ryzen0000001/uts-pl-24-535196-SV-24171}
  \item Pull request \#1: \url{https://github.com/Ryzen0000001/uts-pl-24-535196-SV-24171/pull/1}
  \item Pull request \#2: \url{https://github.com/Ryzen0000001/uts-pl-24-535196-SV-24171/pull/2}
  \item Package backend: \url{https://github.com/users/Ryzen0000001/packages/container/package/uts-pl-24-535196-sv-24171-backend}
  \item Package frontend: \url{https://github.com/users/Ryzen0000001/packages/container/package/uts-pl-24-535196-sv-24171-frontend}
  \item Postman Collection: \url{https://github.com/Ryzen0000001/uts-pl-24-535196-SV-24171/blob/main/postman/buku-servis.postman_collection.json}
\end{itemize}

\end{document}
